\section{Evolución del diseño}

El desarrollo del coche teledirigido no fue lineal, 
sino que evolucionó a lo largo de distintas etapas a 
medida que se evaluaban distintas soluciones.


\subsection{Primera aproximación: Arduino y comunicación por radiofrecuencia}

En una primera fase, se optó por utilizar una placa Arduino 
junto con el módulo de comunicación NRF24L01, con el objetivo de 
implementar un sistema de control remoto basado en radiofrecuencia. 
Esta solución permitía una comunicación directa entre el mando y 
el vehículo, con bajo consumo y relativa simplicidad.

Sin embargo, esta aproximación no cumplía con el objetivo de que 
fuera manejable con el teléfono móvil y requería un Arduino
adicional, además de piezas adicionales como botones, un 
joystick, baterías y el módulo de radiofrecuencia.

\subsection{Transición a ESP32 y servidor web}

Debido a las limitaciones encontradas, se decidió pasar a una 
solución basada en el microcontrolador ESP32, que integra conectividad 
WiFi. Esto permitió implementar un servidor web embebido en el 
propio vehículo, posibilitando el control del coche a través de una 
interfaz accesible desde cualquier dispositivo con navegador.

\subsection{Diseño final}

Finalmente, el sistema se consolidó en torno al ESP32 como núcleo de 
control, integrando la gestión de motores, dirección mediante servomotor 
y comunicación inalámbrica a través de WiFi.
